\documentclass[11pt]{article}
\usepackage[margin=1in]{geometry}
\usepackage{booktabs}
\usepackage{amsmath}
\usepackage{siunitx}
\usepackage{xcolor}
\usepackage{caption}

\sisetup{group-separator={,}, group-minimum-digits=4}

\title{DDM Design Space: Number of Configurations $\mathcal{C}(n, m)$}
\author{CogFormer-2}
\date{}

\begin{document}
\maketitle

\section*{Setup}

A \emph{design configuration} for the DDM is a complete specification of
(i)~which parameters appear in each regressor row of the design matrix, and
(ii)~how many categories each regressor uses.
It is determined by three independently varying factors:

\begin{itemize}
  \item \textbf{Intercept pattern.}  The two dispersion parameters $s_v$ and
    $s_\tau$ are each independently included in or excluded from the
    intercept row, yielding $2^2 = 4$ intercept patterns.  The four
    free parameters $(v, a, z, \tau)$ are always in the intercept; the global
    nuisance parameter $\pi$ is always in the intercept and is never a slope.

  \item \textbf{Slope patterns.}  For each of the $r$ regressors
    ($r \in \{1, \dots, n\}$), any subset of the slope-eligible parameters
    may be assigned.  The slope-eligible set equals the intercept parameters
    minus the global set, so its size depends on the intercept pattern:

    \medskip
    \begin{center}
    \begin{tabular}{lcc}
      \toprule
      Intercept pattern & Slope-eligible params & Patterns / regressor \\
      \midrule
      $\{v, a, z, \tau\}$ (neither $s_v$ nor $s_\tau$)    & 4 & $2^4 = 16$ \\
      $\{v, a, z, \tau, s_v\}$                             & 5 & $2^5 = 32$ \\
      $\{v, a, z, \tau, s_\tau\}$                          & 5 & $2^5 = 32$ \\
      $\{v, a, z, \tau, s_v, s_\tau\}$                     & 6 & $2^6 = 64$ \\
      \bottomrule
    \end{tabular}
    \end{center}

  \item \textbf{Category assignments.}  Each of the $r$ regressors is
    either continuous (1 column) or discrete with $k \in \{2, \dots, m\}$
    categories ($k{-}1$ columns), yielding $m$ choices per regressor.
    Interaction terms are deterministic given the slope patterns (their
    columns are the pairwise intersection of the constituent regressors'
    slopes) and do not add further choices.
\end{itemize}

\section*{Formula}

For $r = 0$ regressors (intercept-only) there are 4 designs.
For $r \geq 1$ regressors the number of designs is
\[
  \bigl(16m\bigr)^r + 2\bigl(32m\bigr)^r + \bigl(64m\bigr)^r,
\]
combining the four intercept patterns with the slope and category
multiplicities.  Summing over all $r \in \{0, 1, \dots, n\}$:

\[
  \mathcal{C}(n, m)
  \;=\;
  4
  \;+\;
  \sum_{r=1}^{n}
  \left[
    (16m)^r + 2(32m)^r + (64m)^r
  \right].
\]

\section*{Results}

Table~\ref{tab:total} reports $\mathcal{C}(n, m)$ for $n, m \in \{2, 3, 4, 5\}$.
Table~\ref{tab:structural} reports the category-agnostic structural count
$\mathcal{S}(n) = \mathcal{C}(n, 1)$, i.e.\ the number of distinct
parameter-assignment structures regardless of how many categories are used.
The four held-out design signatures used in the model-family experiments
are listed in Table~\ref{tab:holdout}.

\begin{table}[h]
\centering
\caption{Total design configurations $\mathcal{C}(n, m)$ for the DDM
  (intercept $+$ up to $n$ regressors, up to $m$ categories each).}
\label{tab:total}
\begin{tabular}{r S[table-format=13.0] S[table-format=13.0] S[table-format=13.0] S[table-format=13.0]}
\toprule
{$n$} & {$m = 2$} & {$m = 3$} & {$m = 4$} & {$m = 5$} \\
\midrule
2 &          25892 &          58036 &         102980 &         160724 \\
3 &        2680100 &        9015988 &       21336644 &       41632724 \\
4 &      305718564 &     1543148212 &     4869952068 &    11879072724 \\
5 &    36846495012 &   279024669364 &  1174174798404 &  3580314272724 \\
\bottomrule
\end{tabular}
\end{table}

\begin{table}[h]
\centering
\caption{Structural design count $\mathcal{S}(n)$ (category-agnostic).
  This equals the number of distinct design-config signatures used by the
  holdout rejection check.}
\label{tab:structural}
\begin{tabular}{r S[table-format=13.0]}
\toprule
{$n$} & {$\mathcal{S}(n)$} \\
\midrule
2 &         6548 \\
3 &       338324 \\
4 &     19278228 \\
5 &   1161177492 \\
\bottomrule
\end{tabular}
\end{table}

\begin{table}[h]
\centering
\caption{The four held-out design configurations for DDM.  These structural
  signatures are excluded from model-family training via the holdout rejection
  loop in \texttt{\_resolve\_design}.  Each represents a negligible fraction
  of $\mathcal{S}(n)$ even at $n = 2$.}
\label{tab:holdout}
\begin{tabular}{lll}
\toprule
Name & Structure & Held-out at $m$ \\
\midrule
\texttt{three\_reg\_interaction}
  & Intercept $+$ 3 regressors $+$ all pairwise interactions
  & $m = 3$ \\
\texttt{two\_factor}
  & Intercept $+$ 2 regressors $+$ interaction, all params on every row
  & $m = 3$ \\
\texttt{asymmetric}
  & Intercept $+$ 2 regressors: drift cluster $(v, z, s_v)$ vs.\ threshold cluster $(a, \tau, s_\tau)$
  & $m = 2$ \\
\texttt{multi\_factor}
  & Intercept $+$ 5 regressors, hierarchical parameter shedding ($6 \to 5 \to 4 \to 3 \to 2$ slopes)
  & $m = 2$ \\
\bottomrule
\end{tabular}
\end{table}

\paragraph{Fraction held out.}
At the default training settings ($n = 5$, $m = 5$), the four holdout
structural signatures represent
$4 / \mathcal{S}(5) \approx 3.4 \times 10^{-9}$
of the structural space and an even smaller fraction of the total design
space $\mathcal{C}(5, 5) \approx 3.58 \times 10^{12}$.

\end{document}
